% ============================================================
% Trabalho 01 — FAED/PPGIA-PUCPR
% Análise Comparativa de Estruturas de Dados
% Formato: IEEE Conference — IEEEtran — duas colunas
% Compilar: pdflatex artigo.tex  (duas passagens)
% ============================================================
\documentclass[conference,a4paper]{IEEEtran}

% --- Pacotes ---
\usepackage[utf8]{inputenc}
\usepackage[T1]{fontenc}
\usepackage[brazil]{babel}
\usepackage{amsmath}
\usepackage{graphicx}
\usepackage{booktabs}
\usepackage{tabularx}
\usepackage{multirow}
\usepackage{siunitx}
\usepackage{url}
\usepackage{hyperref}
\usepackage{microtype}

% Figuras do diretório ../graficos/
\graphicspath{{../graficos/}}

\sisetup{
  group-separator = {.},
  output-decimal-marker = {,},
  table-format = 6.0
}

% ============================================================
\begin{document}

\title{Análise Comparativa de Estruturas de Dados:\\
Arrays, Árvores de Busca Binária e Tabelas Hash}

\author{%
  \IEEEauthorblockN{Dantas, F.~S.\textsuperscript{1} \and
                    Correia, R.\textsuperscript{1}}
  \IEEEauthorblockA{%
    \textsuperscript{1}Programa de Pós-Graduação em Informática Aplicada (PPGIA)\\
    Pontifícia Universidade Católica do Paraná (PUCPR)\\
    Curitiba, Brasil\\
    \{fsd-dantas, rafacs2002\}@pucpr.edu.br}
}

\maketitle

% ============================================================
\begin{abstract}
Este artigo apresenta uma análise comparativa experimental entre três
categorias de estruturas de dados—arrays lineares, árvores de busca
binária (BST) com e sem balanceamento AVL, e tabelas hash com
encadeamento externo—avaliando operações de inserção e busca em volumes
de registros $N \in \{10\,000;\,50\,000;\,100\,000\}$. As métricas
coletadas incluem número de iterações, tempo de execução, consumo de
memória via \texttt{tracemalloc} e variação de RSS via \texttt{psutil},
com pelo menos cinco rodadas independentes por cenário. Os resultados
confirmam as complexidades assintóticas teóricas e revelam nuances
práticas relevantes: o custo do balanceamento AVL durante a inserção,
o impacto crítico do fator de carga $\alpha$ na tabela hash, e a
degradação da função de dobramento com chaves numéricas de 9 dígitos.
\end{abstract}

\begin{IEEEkeywords}
estruturas de dados; busca binária; árvore AVL; tabela hash; análise
de desempenho; complexidade assintótica
\end{IEEEkeywords}

% ============================================================
\section{Introdução}

A escolha da estrutura de dados adequada é uma decisão fundamental em
engenharia de software, com impacto direto na eficiência computacional
das aplicações. Embora as complexidades assintóticas de arrays, árvores
de busca binária e tabelas hash sejam amplamente conhecidas na
literatura~\cite{cormen2022}, a tradução dessas garantias teóricas para
comportamento empírico depende de fatores como distribuição dos dados,
parametrizações internas e características do hardware.

Este trabalho tem como objetivo implementar, do zero em Python~3.13, as
três estruturas mencionadas e submetê-las a um protocolo experimental
controlado. A contribuição principal reside na análise conjunta de
múltiplas métricas—iterações, tempo, memória—sob diferentes volumes e
parametrizações, relacionando os resultados à teoria de complexidade
assintótica e discutindo os \textit{trade-offs} de cada abordagem.

O restante do artigo está organizado como segue: a Seção~\ref{sec:teoria}
apresenta o referencial teórico; a Seção~\ref{sec:metodologia} descreve
a metodologia experimental; a Seção~\ref{sec:implementacao} detalha as
implementações; a Seção~\ref{sec:resultados} analisa os resultados; e a
Seção~\ref{sec:conclusao} apresenta as conclusões.

% ============================================================
\section{Referencial Teórico}
\label{sec:teoria}

\subsection{Array Linear}

Um array linear armazena elementos em posições contíguas de memória.
A inserção ao final é $O(1)$ amortizado em estruturas dinâmicas. A
busca linear percorre o array até encontrar o elemento—$O(n)$ no pior
caso—enquanto a busca binária, que exige ordem prévia, realiza $O(\log n)$
comparações por divisão e conquista~\cite{knuth1998}.

\subsection{Árvore de Busca Binária}

Uma BST mantém a propriedade $\text{esq} < \text{raiz} < \text{dir}$ em
cada nó, garantindo $O(\log n)$ para inserção e busca em árvores
balanceadas. No entanto, sem mecanismo de balanceamento, inserções em
ordem crescente ou decrescente degeneram a árvore em lista encadeada,
elevando o custo ao pior caso $O(n)$~\cite{cormen2022}.

A árvore AVL~\cite{adelson1962} restaura o balanceamento após cada
inserção por meio de rotações simples (LL, RR) e duplas (LR, RL),
mantendo o fator de balanço $|h_e - h_d| \leq 1$ em todo nó. Isso
garante altura $h \leq 1{,}44\,\log_2(N+2)$, ou seja, $O(\log n)$
estrito para todas as operações.

\subsection{Tabela Hash}

Uma tabela hash de tamanho $M$ mapeia cada chave $k$ a um índice
$h(k) \in [0, M)$. O fator de carga $\alpha = n/M$ quantifica a
densidade da tabela; com encadeamento externo, o custo esperado de
busca é $O(1 + \alpha)$~\cite{knuth1998}. Três funções foram avaliadas:

\begin{itemize}
  \item \textbf{Divisão}: $h(k) = k \bmod M$. Simples e eficiente;
        sensível à escolha de $M$ (evitar potências de~2).
  \item \textbf{Multiplicação} (Knuth): $h(k) = \lfloor M \cdot
        \{k \cdot A\} \rfloor$, com $A = (\sqrt{5}-1)/2 \approx 0{,}618$.
        Menos sensível ao valor de $M$.
  \item \textbf{Dobramento}: soma de blocos de 3 dígitos da chave,
        seguida de $\bmod M$. Indicada para chaves numéricas longas.
\end{itemize}

% ============================================================
\section{Metodologia Experimental}
\label{sec:metodologia}

\subsection{Geração dos Dados}

Os registros fictícios simulam cadastros de funcionários, com campos:
\textit{matrícula} (inteiro de 9 dígitos, chave primária única),
\textit{nome}, \textit{salário} e \textit{código do setor}. A geração
utiliza semente fixa (\texttt{seed=42}), garantindo que todas as
estruturas recebam exatamente os mesmos dados em todas as rodadas.

\subsection{Protocolo de Medição}

Para cada combinação de estrutura, volume $N$ e parametrização,
executaram-se $r = 5$ rodadas independentes. A Tabela~\ref{tab:cenarios}
resume os cenários avaliados.

\begin{table}[htbp]
  \caption{Cenários de Teste}
  \label{tab:cenarios}
  \centering
  \begin{tabular}{lll}
    \toprule
    Estrutura & Variações & $N$ \\
    \midrule
    Array (linear/binária) & — & $10\,k,\,50\,k,\,100\,k$ \\
    BST / AVL              & sem / com balanceamento & $10\,k,\,50\,k,\,100\,k$ \\
    Hash (encadeamento)    & $M \in \{100,1000,5000\}$ & $10\,k,\,50\,k,\,100\,k$ \\
                           & 3 funções hash & \\
    \bottomrule
  \end{tabular}
\end{table}

As métricas coletadas em cada rodada foram:

\begin{itemize}
  \item \textbf{Iterações}: comparações realizadas na operação;
  \item \textbf{Tempo}: \texttt{time.perf\_counter()} em segundos;
  \item \textbf{Memória Python}: pico de alocação via \texttt{tracemalloc};
  \item \textbf{RSS delta}: variação de memória residente via \texttt{psutil}.
\end{itemize}

Os experimentos foram executados em ambiente local controlado
(Windows~11, Python~3.13.12, processador x86-64), sem carga paralela
significativa, conforme recomendado pelo enunciado.

% ============================================================
\section{Implementação}
\label{sec:implementacao}

Todas as estruturas foram implementadas sem utilizar bibliotecas que já
as forneçam. Apenas \texttt{tracemalloc}, \texttt{psutil},
\texttt{time} e \texttt{csv} foram empregados para medição e I/O.

\subsection{Array Linear}

A classe \texttt{LinearArray} encapsula uma lista Python com operações
de inserção (\texttt{append}, $O(1)$ amortizado), busca linear
($O(n)$) e busca binária ($O(\log n)$). Para a busca binária,
implementou-se QuickSort iterativo (com pilha explícita) para evitar
\textit{stack overflow} em $N = 100\,000$.

\subsection{BST e AVL}

A \texttt{BST} usa inserção iterativa com contador de comparações.
A \texttt{AVL} adota inserção recursiva com atualização de altura e
aplicação das quatro rotações de rebalanceamento após cada inserção.
Ambas expõem \texttt{inserir()}, \texttt{buscar()} e \texttt{altura()}.

\subsection{Tabela Hash}

A classe \texttt{HashTable(m, funcao)} implementa encadeamento externo
com lista encadeada por balde. O método \texttt{estatisticas\_cadeias()}
retorna comprimento máximo e médio das cadeias, baldes vazios e fator de
carga—métricas essenciais para avaliar a qualidade da função hash.

% ============================================================
\section{Resultados e Análise}
\label{sec:resultados}

\subsection{Busca: Iterações vs.\ Volume}

A Fig.~\ref{fig:busca_iter} ilustra o crescimento de iterações de busca
em função de $N$.

\begin{figure}[htbp]
  \centering
  \includegraphics[width=\columnwidth]{fig1_busca_iteracoes.png}
  \caption{Iterações de busca vs.\ $N$ (escala logarítmica).}
  \label{fig:busca_iter}
\end{figure}

A busca linear cresce linearmente: para $N=100\,000$, o alvo na
posição central exige $\approx 50\,001$ comparações. A busca binária,
BST e AVL permanecem em escala logarítmica. Notavelmente, a hash com
$M=5\,000$ (divisão) alcança apenas 11 iterações para $N=100\,000$,
confirmando $O(1+\alpha)$ com $\alpha = 20$.

\subsection{Inserção: Tempo vs.\ Volume}

A Fig.~\ref{fig:insercao_tempo} compara o tempo de inserção.

\begin{figure}[htbp]
  \centering
  \includegraphics[width=\columnwidth]{fig2_insercao_tempo.png}
  \caption{Tempo de inserção vs.\ $N$.}
  \label{fig:insercao_tempo}
\end{figure}

O array exibe o menor tempo de inserção em todos os volumes—confirmando
$O(1)$ amortizado. A AVL é consistentemente mais lenta que a BST
($\approx\!2{,}5\times$ para $N=100\,000$) em razão do custo das
rotações. A hash com $M=100$ é a mais lenta: com $\alpha=1\,000$, cada
inserção percorre uma cadeia longa antes de encontrar a posição.

\subsection{Impacto do Fator de Carga na Hash}

A Tabela~\ref{tab:hash_busca} e a Fig.~\ref{fig:hash_m} mostram o
impacto de $M$ e da função hash nas iterações de busca para $N=100\,000$.

\begin{table}[htbp]
  \caption{Tabela Hash — Busca (iterações, $N=100\,000$, média de 5 rodadas)}
  \label{tab:hash_busca}
  \centering
  \begin{tabular}{rrrr}
    \toprule
    $M$ & $\alpha$ & Divisão & Dobramento \\
    \midrule
    100   & 1000 & 520 & 481 \\
    1\,000 & 100  &  44 &  52 \\
    5\,000 &  20  &  11 &  34 \\
    \bottomrule
  \end{tabular}
\end{table}

\begin{figure}[htbp]
  \centering
  \includegraphics[width=\columnwidth]{fig3_hash_impacto_m.png}
  \caption{Impacto de $M$ nas iterações de busca ($N=100\,000$).}
  \label{fig:hash_m}
\end{figure}

Com $M=100$ e $N=100\,000$, o fator de carga $\alpha=1\,000$ faz a
hash degenerar em $O(n)$—comportamento indistinguível da busca linear.
Com $M=5\,000$ ($\alpha=20$), a hash (divisão) supera todas as demais
estruturas em velocidade de busca.

\subsection{Comparação das Funções Hash}

A Fig.~\ref{fig:hash_fn} revela uma divergência significativa do
método de dobramento em $M=5\,000$.

\begin{figure}[htbp]
  \centering
  \includegraphics[width=\columnwidth]{fig4_hash_funcoes.png}
  \caption{Iterações de inserção por função hash e $M$ ($N=100\,000$).}
  \label{fig:hash_fn}
\end{figure}

O dobramento gerou $2{,}8\times$ mais iterações de inserção que a
divisão em $M=5\,000$ ($2\,841\,623$ vs.\ $998\,494$). A causa é o
\textit{clustering}: matrículas de 9 dígitos no intervalo
$[10^8, 10^9)$ produzem somas de blocos concentradas em faixas
estreitas, reduzindo a efetividade da dispersão. Este resultado
evidencia que a adequação de uma função hash depende da distribuição
das chaves, não apenas da sua formulação matemática.

\subsection{BST vs.\ AVL: Altura e Balanceamento}

A Fig.~\ref{fig:altura} compara a altura das árvores com a referência
teórica $\lfloor \log_2 N \rfloor$.

\begin{figure}[htbp]
  \centering
  \includegraphics[width=\columnwidth]{fig5_bst_avl_altura.png}
  \caption{Altura das árvores vs.\ $N$, com referência $\log_2 N$.}
  \label{fig:altura}
\end{figure}

\begin{table}[htbp]
  \caption{Altura das Árvores por Volume}
  \label{tab:altura}
  \centering
  \begin{tabular}{rrrr}
    \toprule
    $N$ & BST & AVL & $\lfloor\log_2 N\rfloor$ \\
    \midrule
    10\,000  & 29 & 15 & 13 \\
    50\,000  & 35 & 18 & 15 \\
    100\,000 & 38 & 19 & 16 \\
    \bottomrule
  \end{tabular}
\end{table}

A AVL mantém-se consistentemente próxima do ideal ($h_{\text{AVL}}
\leq 1{,}44\,\log_2 N$), enquanto a BST com dados aleatórios já
apresenta altura $\approx\!2\times$ a ideal. Em dados ordenados, a BST
degeneraria completamente—cenário não testado aqui, mas amplamente
documentado~\cite{cormen2022}.

\subsection{Análise de Memória}

A Tabela~\ref{tab:memoria} compara o consumo de memória durante a inserção.

\begin{table}[htbp]
  \caption{Pico de Memória Python na Inserção ($N=100\,000$)}
  \label{tab:memoria}
  \centering
  \begin{tabular}{lS[table-format=2.2]}
    \toprule
    Estrutura & {Pico (MB)} \\
    \midrule
    Array linear       & 0{,}76 \\
    BST                & 6{,}10 \\
    AVL                & 6{,}87 \\
    Hash ($M=5\,000$)  & 5{,}38 \\
    Hash ($M=1\,000$)  & 5{,}35 \\
    Hash ($M=100$)     & 5{,}34 \\
    \bottomrule
  \end{tabular}
\end{table}

O array consome $\approx\!8\times$ menos memória que as árvores, pois
não aloca objetos por nó. A AVL exige levemente mais que a BST pelo
campo de altura adicional em cada nó. O tamanho $M$ da hash impacta
pouco a memória—o custo dominante são os nós das cadeias, não a tabela.

% ============================================================
\section{Conclusão}
\label{sec:conclusao}

Este trabalho realizou uma análise comparativa experimental de arrays
lineares, árvores BST/AVL e tabelas hash, confirmando as complexidades
assintóticas teóricas e revelando aspectos práticos relevantes.

Os principais achados são:

\begin{enumerate}
  \item \textbf{Array linear}: menor custo de memória e inserção; busca
        linear $O(n)$ inviável para grandes volumes; busca binária
        $O(\log n)$ competitiva mas exige ordenação prévia $O(n\log n)$.

  \item \textbf{BST}: equilibrada para dados aleatórios; vulnerável à
        degeneração com dados ordenados. Inserção mais rápida que AVL.

  \item \textbf{AVL}: garantia $O(\log n)$ estrita; altura próxima do
        ideal teórico; custo de inserção $\approx\!2{,}5\times$ maior
        que BST pelo overhead das rotações.

  \item \textbf{Hash ($M$ adequado)}: melhor desempenho de busca quando
        $\alpha$ é pequeno ($M\!=\!5\,000$, $\alpha\!=\!20$: 11
        iterações para $N\!=\!100\,000$). Com $M$ pequeno ($\alpha\!
        =\!1\,000$), degenera em $O(n)$.

  \item \textbf{Função hash}: divisão e multiplicação distribuem bem
        matrículas de 9 dígitos; dobramento apresenta \textit{clustering}
        que eleva as colisões em 2,8× para $M=5\,000$.
\end{enumerate}

A escolha da estrutura deve considerar o perfil da aplicação: arrays
para leitura sequencial com memória limitada; AVL para garantias de
busca sob qualquer ordenação de inserção; hash com $M\!\gg\!n$ para
acesso mais próximo de $O(1)$ quando memória não é restrição crítica.

% ============================================================
\begin{thebibliography}{9}

\bibitem{cormen2022}
T.~H. Cormen, C.~E. Leiserson, R.~L. Rivest, e C.~Stein,
\textit{Introduction to Algorithms}, 4ª~ed.
Cambridge, MA: MIT Press, 2022.

\bibitem{knuth1998}
D.~E. Knuth,
\textit{The Art of Computer Programming, Vol.~3: Sorting and Searching},
2ª~ed.
Reading, MA: Addison-Wesley, 1998.

\bibitem{adelson1962}
G.~M. Adelson-Velsky e E.~M. Landis,
``An algorithm for the organization of information,''
\textit{Proceedings of the USSR Academy of Sciences},
vol.~146, pp.~263--266, 1962.

\bibitem{python2024}
Python Software Foundation,
\textit{Python 3.13 Documentation}.
Disponível em: \url{https://docs.python.org/3.13/}.
Acesso em: mar. 2026.

\bibitem{sedgewick2011}
R.~Sedgewick e K.~Wayne,
\textit{Algorithms}, 4ª~ed.
Upper Saddle River, NJ: Addison-Wesley, 2011.

\end{thebibliography}

\end{document}
